% ============================================================
% Section 25: 一般幂律势的晶体结合能与体积弹性模量
% ============================================================
\section{固体理论：一般幂律势的晶体结合能与体积弹性模量}
\label{sec:power-law-potential}

\begin{modelbox}
设某一晶体中相邻原子间相互作用势能为
\[
  u(r) = -\frac{\alpha}{r^m} + \frac{\beta}{r^n} \qquad (n > m > 0).
\]
\begin{enumerate}[nosep]
  \item 若只计及最近邻原子的相互作用势能，证明其体积弹性模量为
    \[
      K = \frac{|U_0|\, mn}{9V_0},
    \]
    其中 $U_0$ 为晶体在平衡时的总相互作用势能，$V_0$ 为平衡体积。
  \item 设平衡时两原子间距 $r_0 = 3\,\text{\mathring{\mathrm{A}}}$，每个原子的离解能为 $4\,\text{eV}$，$m=2$, $n=10$，计算 $\alpha$ 及 $\beta$ 值。
\end{enumerate}
\end{modelbox}

\subsection{概念铺垫：体积弹性模量的定义}

\begin{modelbox}{体积模量 $K$ 的两个等价定义——从何而来？}

\textbf{定义 1：从压缩实验出发}

想象你把一个固体块放进液压机，均匀增加外部压强 $P$。压强增加 $dP$ 时，体积被压缩了 $dV$（注意 $dV<0$）。

你想定义一个量来衡量``这块材料有多难被压扁''。直觉上：
\begin{itemize}[nosep]
  \item 若压强只增加一点点，体积就大幅缩小$\Rightarrow$材料很软；
  \item 若压强增加很多，体积几乎不变$\Rightarrow$材料很硬。
\end{itemize}

因此，体积模量应该正比于``压强增量''与``相对体积变化''的比值：
\[
  K = \frac{\text{压强增量}}{\text{相对体积压缩}} = \frac{dP}{-dV/V} = -V\,\frac{dP}{dV}.
\]
\textbf{负号的必要性}：压强增加（$dP>0$）时体积减小（$dV<0$），负号保证 $K>0$。

\textbf{类比杨氏模量}：$Y = \dfrac{F/A}{\Delta L/L_0} = \dfrac{\text{应力}}{\text{应变}}$。体积模量完全类似，只是把``线应变''换成了``体积应变''$-(dV/V)$。

\vspace{0.5em}
\textbf{定义 2：从能量曲率推导}

在热力学中，压强是能量对体积的负一阶导（体积增大时系统对外做功，内能减小）：
\[
  P = -\frac{dU}{dV}.
\]

对压强再求一次体积导数：
\[
  \frac{dP}{dV} = -\frac{d^2U}{dV^2}.
\]

代入定义 1：
\[
  K = -V\,\frac{dP}{dV} = -V\left(-\frac{d^2U}{dV^2}\right) = \boxed{V\,\frac{d^2U}{dV^2}}.
\]

\textbf{物理图像}：$d^2U/dV^2$ 是能量$-$体积曲线在平衡点的\textbf{曲率}。曲率越大，意味着体积稍有偏离平衡，能量就急剧升高$\Rightarrow$材料强烈``抗拒''形变，即 $K$ 很大。钻石的能量曲线曲率极大，故体积模量 $K\sim 440\,\text{GPa}$；橡胶的曲率很小，$K$ 也极小。

\textbf{为什么乘以 $V$ 而不是别的？} 做一个检验：把两个完全相同的系统拼在一起，总体积 $2V$，总能量 $2U$。
\[
  \frac{d^2(2U)}{d(2V)^2} = \frac{2U''}{4} = \frac{U''}{2}.
\]
若定义 $K$ 时不乘 $V$，则合并后 $K$ 变成原来的一半——这不对，因为材料属性不应随系统大小改变。乘以 $V$ 后：
\[
  K_{\text{合并}} = (2V)\cdot\frac{U''}{2} = VU'' = K_{\text{原}}.
\]
因此 $V$ 因子确保 $K$ 是\textbf{强度量}（与系统大小无关的材料常数）。
\end{modelbox}

\begin{insightbox}{为什么体积模量与晶体结构无关？}
晶体体积 $V$ 与最近邻间距 $r$ 的关系为 $V = NAr^3$，其中 $A$ 是由晶体结构决定的几何常数：
\begin{center}
\renewcommand{\arraystretch}{1.4}
\begin{tabular}{|l|c|c|}
\hline
\textbf{结构} & \textbf{几何常数} $A$ & \textbf{每个原子体积} \\
\hline
简单立方（sc） & $1$ & $r^3$ \\
面心立方（fcc） & $\frac{1}{\sqrt{2}}$ & $\frac{r^3}{\sqrt{2}}$ \\
体心立方（bcc） & $\frac{4}{3\sqrt{3}}$ & $\frac{3\sqrt{3}}{4}r^3$ \\
\hline
\end{tabular}
\end{center}

在最终公式中，$A$ 会被消掉，因此体积模量与晶体结构的细节无关——这是幂律势的一个美妙性质。
\end{insightbox}

\subsection{体积弹性模量的证明}

\subsubsection{总相互作用能与平衡条件}

\begin{modelbox}{$N/2$ 还是 $Nz/2$？——计数规则的澄清}
在固体物理中，如果只计最近邻相互作用，总能量有两种等价写法，取决于你如何定义 $u(r)$：

\textbf{写法 1（严格计数，每对原子）}

$u(r)$ 是\textbf{一对}最近邻原子的相互作用能。每个原子有 $z$ 个最近邻，总键数（不重复计数）为 $Nz/2$：
\[
  U(r) = \frac{Nz}{2}\,u(r) = \frac{Nz}{2}\left(-\frac{\alpha}{r^m} + \frac{\beta}{r^n}\right).
\]

\textbf{写法 2（本题采用的约定）}

把 $z$ 吸收进势能参数中，即令 $\tilde{\alpha} = z\alpha$, $\tilde{\beta} = z\beta$。此时等效的单键势已经隐含了配位数：
\[
  U(r) = \frac{N}{2}\left(-\frac{\tilde{\alpha}}{r^m} + \frac{\tilde{\beta}}{r^n}\right).
\]
本题以及多数教材在推导 $K = |U_0|mn/(9V_0)$ 时采用写法 2，因为最终结果与 $z$ 无关，这样公式更简洁。

\textbf{关键原则}：$\alpha$ 和 $\beta$ 的数值\textbf{依赖于你采用的约定}。如果你用写法 1，算出的 $\alpha$ 会比写法 2 小 $z$ 倍。但物理结果（$K$, $r_0$）不变。
\end{modelbox}

只计及最近邻相互作用，总势能为
\[
  U(r) = \frac{N}{2}\left(-\frac{\alpha}{r^m} + \frac{\beta}{r^n}\right). \tag{1}
\]

平衡条件 $\dd U/\dd r|_{r_0} = 0$：
\[
  \frac{N}{2}\left(\frac{m\alpha}{r_0^{m+1}} - \frac{n\beta}{r_0^{n+1}}\right) = 0
  \quad\Longrightarrow\quad
  \frac{m\alpha}{r_0^m} = \frac{n\beta}{r_0^n}. \tag{2}
\]

由此解出平衡间距：
\[
  r_0 = \left(\frac{n\beta}{m\alpha}\right)^{\!\frac{1}{n-m}}. \tag{3}
\]

将 (2) 代入 (1) 得平衡时的总势能：
\begin{align*}
  U_0 &= \frac{N}{2}\left(-\frac{\alpha}{r_0^m} + \frac{m}{n}\frac{\alpha}{r_0^m}\right)
      = -\frac{N}{2}\frac{\alpha}{r_0^m}\left(1 - \frac{m}{n}\right) \\
      &= -\frac{N\alpha}{2r_0^m}\cdot\frac{n-m}{n}. \tag{4}
\end{align*}

\begin{tipbox}{符号约定}
$U_0 < 0$（结合能），因此 $|U_0| = -U_0$。式 (4) 给出
\[
  \frac{\alpha}{r_0^m} = \frac{2|U_0|}{N}\cdot\frac{n}{n-m}. \tag{4'}
\]
\end{tipbox}

\subsubsection{从 $r$ 到 $V$ 的变量转换}

晶体体积与原子间距的关系：
\[
  V = NAr^3 \quad\Longrightarrow\quad \frac{\dd V}{\dd r} = 3NAr^2,
  \quad \frac{\dd r}{\dd V} = \frac{1}{3NAr^2}. \tag{5}
\]

由链式法则，注意在平衡点 $r_0$ 处一阶导 $\dd U/\dd r = 0$：
\[
  \frac{\dd^2 U}{\dd V^2}
  = \frac{\dd}{\dd V}\!\left(\frac{\dd U}{\dd r}\frac{\dd r}{\dd V}\right)
  = \frac{\dd^2 U}{\dd r^2}\left(\frac{\dd r}{\dd V}\right)^{\!2}
    + \underbrace{\frac{\dd U}{\dd r}}_{=0}\frac{\dd^2 r}{\dd V^2}.
\]

因此
\[
  K = V_0\left.\frac{\dd^2 U}{\dd V^2}\right|_{V_0}
    = V_0\left(\frac{1}{3NAr_0^2}\right)^{\!2}\left.\frac{\dd^2 U}{\dd r^2}\right|_{r_0}. \tag{6}
\]

\subsubsection{计算二阶导数}

\[
  \frac{\dd^2 U}{\dd r^2}
  = \frac{N}{2}\left[-\frac{m(m+1)\alpha}{r^{m+2}} + \frac{n(n+1)\beta}{r^{n+2}}\right].
\]

在 $r_0$ 处利用 (2) 式替换 $\beta/r_0^n = m\alpha/(nr_0^m)$：
\begin{align*}
  \left.\frac{\dd^2 U}{\dd r^2}\right|_{r_0}
  &= \frac{N\alpha}{2r_0^{m+2}}\left[-m(m+1) + n(n+1)\cdot\frac{m}{n}\right] \\
  &= \frac{N\alpha}{2r_0^{m+2}}\left[-m(m+1) + m(n+1)\right] \\
  &= \frac{N\alpha}{2r_0^{m+2}}\,m(n-m). \tag{7}
\end{align*}

\subsubsection{代入化简}

将 (5) 和 (7) 代入 (6)：
\begin{align*}
  K &= (NAr_0^3)\cdot\frac{1}{9N^2A^2r_0^4}\cdot\frac{N\alpha\,m(n-m)}{2r_0^{m+2}} \\
    &= \frac{m(n-m)}{18Ar_0^{m+3}}\cdot\alpha.
\end{align*}

利用 (4') 式 $\alpha/r_0^m = 2|U_0|n/[N(n-m)]$：
\begin{align*}
  K &= \frac{m(n-m)}{18Ar_0^3}\cdot\frac{2|U_0|}{N}\cdot\frac{n}{n-m} \\
    &= \frac{2mn|U_0|}{18NAr_0^3}
     = \frac{|U_0|\,mn}{9\underbrace{NAr_0^3}_{=V_0}}.
\end{align*}

故
\[
  \boxed{ K = \frac{|U_0|\,mn}{9V_0} }. \tag{证毕}
\]

\begin{tipbox}{考场速记}
对幂律势 $-\alpha/r^m + \beta/r^n$，体积模量与结合能的关系被 $mn/9$ 锁定：
\begin{center}
\renewcommand{\arraystretch}{1.4}
\begin{tabular}{|c|c|c|l|}
\hline
\textbf{体系} & $m$ & $n$ & $\bm{K\cdot 9V_0/|U_0|}$ \\
\hline
离子晶体（NaCl） & 1 & 9 & 9 \\
惰性气体（LJ） & 6 & 12 & 72 \\
本题 & 2 & 10 & 20 \\
\hline
\end{tabular}
\end{center}
\end{tipbox}

\subsection{参数 $\alpha$ 与 $\beta$ 的计算}

已知：$r_0 = 3\,\text{\mathring{\mathrm{A}}}$，每个原子的离解能 $|U_0|/N = 4\,\text{eV}$，$m=2$, $n=10$。

\textbf{离解能}的物理意义：将晶体完全拆散成孤立原子所需的能量，即 $|u_0| = |U_0|/N$。

\subsubsection{求 $\alpha$}

由 (4) 式：
\[
  \frac{|U_0|}{N} = \frac{1}{2}\frac{\alpha}{r_0^m}\cdot\frac{n-m}{n}.
\]

代入数据：
\[
  4 = \frac{1}{2}\cdot\frac{\alpha}{3^2}\cdot\frac{10-2}{10}
  = \frac{1}{2}\cdot\frac{\alpha}{9}\cdot\frac{8}{10}
  = \frac{2\alpha}{45}.
\]

解得
\[
  \boxed{ \alpha = 90\,\text{eV\cdot\mathring{\mathrm{A}}}^2 }.
\]

\subsubsection{求 $\beta$}

由平衡条件 (3)：
\[
  r_0^{n-m} = \frac{n\beta}{m\alpha}
  \quad\Longrightarrow\quad
  \beta = \frac{m\alpha}{n}\,r_0^{n-m}.
\]

代入 $m=2$, $n=10$, $\alpha=90$, $r_0=3$：
\[
  \beta = \frac{2\times 90}{10}\times 3^{10-2}
       = 18\times 3^8
       = 18\times 6561
       = 118098.
\]

故
\[
  \boxed{ \beta \approx 1.18\times 10^5\,\text{eV\cdot\mathring{\mathrm{A}}}^{10} }.
\]

\begin{tipbox}{量纲检验}
势函数 $u(r)$ 的量纲是能量。因此：
\begin{itemize}[nosep]
  \item $[\alpha/r^m] = \text{能量} \Rightarrow [\alpha] = \text{能量}\cdot\text{长度}^m$。
  \item 本题 $m=2$：$[\alpha] = \text{eV\cdot\mathring{\mathrm{A}}}^2$。正确。
  \item 本题 $n=10$：$[\beta] = \text{eV\cdot\mathring{\mathrm{A}}}^{10}$。正确。
\end{itemize}
\end{tipbox}

\subsection{物理图像与拓展}

\subsubsection{幂律势中 $m$ 与 $n$ 的物理来源}

\begin{center}
\renewcommand{\arraystretch}{1.5}
\begin{tabular}{|l|c|c|l|}
\hline
\textbf{相互作用类型} & $m$ & $n$ & \textbf{物理来源} \\
\hline
库仑（离子晶体） & 1 & 9 & 吸引：$\pm e^2/(4\pi\varepsilon_0 r)$；排斥：电子云交叠 \\
范德华（惰性气体） & 6 & 12 & 吸引：诱导偶极 $-C_6/r^6$；排斥：Pauli 不相容 \\
偶极-偶极 & 3 & $>3$ & Keesom 相互作用 \\
\hline
\end{tabular}
\end{center}

排斥项的幂次 $n$ 总是大于吸引项的幂次 $m$，这是因为排斥力（电子云交叠、Pauli 原理）是\textbf{短程}的，随距离下降更快。

\subsubsection{体积模量的数量级}

由 $K = |U_0|mn/(9V_0)$，结合能 $|U_0|/N \sim \text{eV/atom}$，原子体积 $V_0/N \sim \text{\mathring{\mathrm{A}}}^3$，得
\[
  K \sim \frac{(\text{eV})(1)(10)}{9(\text{\mathring{\mathrm{A}}}^3)} \sim 10\,\text{eV/\mathring{\mathrm{A}}}^3 \sim 10^{11}\,\text{Pa} = 10^2\,\text{GPa}.
\]
这与典型固体的体积模量（金属 $10$–$100\,\text{GPa}$，钻石 $440\,\text{GPa}$）同一数量级。
